\documentclass[12pt,a4paper]{article}

\usepackage[T1]{fontenc}
\usepackage[margin=1in]{geometry}
\usepackage{fancyhdr}
\usepackage{enumitem}
\usepackage{setspace}

\setlength{\parindent}{0pt}
\setlength{\parskip}{8pt}
\onehalfspacing

% Footer
\pagestyle{fancy}
\fancyhf{}
\renewcommand{\headrulewidth}{0pt}
\fancyfoot[L]{Blurguard Project Journal}
\fancyfoot[R]{\thepage}

\begin{document}

\begin{center}
    {\Large\bfseries PROJECT JOURNAL}\\[10pt]
    {\large\bfseries BlurGuard -- Real-Time Screen-Share Sensitive Data Censor}\\[10pt]
    {\bfseries Journal : October 2026}
\end{center}

\vspace{6pt}

\textbf{Team Members:}
\begin{itemize}[topsep=2pt,itemsep=0pt]
    \item Lakshita
    \item Rhythm
    \item Yashjot
\end{itemize}

\section*{Work Completed}
This week we were told to create class and component diagram for our blurgaurd application. We divided the work between us, I took the responsibility for class diagram and started with writing each class attributes and functionalities, then joined the those classes according to appropriate relations they share. The class diagram was further used for referencing while making component diagram. I have also made the component diagram.

\section*{Outcome}
The class and component diagram was completed, which can be used to understand project component wise during development phase.

\section*{Next Planned Work}
To create DFD (level 0,level 1,level 2) and add functionality of more modules and make it work for real time data detection and masking on application level.

\end{document}
